\documentclass[11pt, a4paper]{article}

\usepackage[a4paper, top=2.5cm, bottom=2.5cm, left=2cm, right=2cm]{geometry}
\usepackage{amsmath}
\usepackage{amssymb}
\usepackage[colorlinks=true, linkcolor=blue, urlcolor=blue, citecolor=blue]{hyperref}

\title{Theory of Everything - Volume 24 \\ \Large Non-Locality \& Bell's Theorem}
\author{Omega Protocol}
\date{\today}

\begin{document}

\maketitle

\section*{Layman's Summary}
Einstein hated the idea that measuring a particle here could instantly affect a particle on the other side of the galaxy. He called it "spooky." Volume 24 proves that the universe is, indeed, spooky. The universe is fundamentally non-local; distance is an illusion created by the limitations of how we interact with the world.

\section{Violation of Bell Inequalities (Axiom 27)}
John Bell proved that no "local hidden variables"—secret classical rules running behind the scenes—can explain quantum mechanics. We establish this as a core axiom. Local realism is strictly prohibited by the non-commutative structure of the universe's fundamental algebra. Things do not have definite properties until they interact.

\section{Contextuality (Theorem)}
The Kochen-Specker theorem suggests that asking a particle "how fast are you going?" gives a different answer depending on what other questions you ask at the same time. We prove that contextuality is absolute: the outcome of any quantum measurement is completely dependent on the entire geometric and informational context of the experimental setup.

\section{Action at a Distance (Theorem)}
If things can instantly influence each other across the galaxy, doesn't that break the speed of light? We prove that causality is perfectly preserved. Apparent non-local influences propagate instantly because they travel through the "Archive Mode" (the underlying, non-spatial data structure of the universe), not through the 3D space we see.

\end{document}
